\chapter{مبرهنة الأعداد الأولية في المتتاليات الحسابية}

\section{نطاق الفصل وحالته}

\noindent\textbf{حالة الفصل:} \textenglish{VERIFIED}. يثبت هذا الفصل، لكل
ترديد ثابت \(q\ge2\) ولكل فئة مختزلة \(a\pmod q\)، الصيغ النوعية
\[
\psi(x;q,a)\sim\frac{x}{\varphi(q)},
\qquad
\vartheta(x;q,a)\sim\frac{x}{\varphi(q)},
\qquad
\pi(x;q,a)\sim\frac{x}{\varphi(q)\log x}.
\]

بدأ التأليف بعد إغلاق بوابة
\textenglish{EVIDENCE-FIRST / PRE-AUTHORING}. وتوجد مواد التدقيق في:
\begin{itemize}
  \item \texttt{research/literature-reviews/chapter-10-prime-number-theorem-}
  \texttt{arithmetic-progressions-evidence.md}.
  \item \texttt{research/literature-reviews/chapter-10-prime-number-theorem-}
  \texttt{arithmetic-progressions-proof-map.md}.
  \item \texttt{docs/CHAPTER\_10\_PRE\_AUTHORING\_AUDIT\_2026-07-20.md}.
  \item \texttt{docs/CHAPTER\_10\_LOGIC\_AUDIT\_2026-07-20.md}.
  \item \texttt{docs/CHAPTER\_10\_BIBLIOGRAPHIC\_VERIFICATION\_2026-07-20.md}.
  \item \texttt{docs/CHAPTER\_10\_AUTHORING\_AUDIT\_2026-07-20.md}.
\end{itemize}

تعني حالة \textenglish{VERIFIED} أن النتائج الداخلية السبع اجتازت تدقيق
الإشارات وعدم الدور، والتحقق المرجعي، وتدقيق ما بعد التأليف. ولا تعني هذه
الحالة \textenglish{REVIEWED}؛ فالمراجعة الثانية المستقلة ما تزال مطلوبة.

المسار تاوبيري نوعي: نفكك سلسلة فون مانغولت في الفئة بواسطة شخصيات
ديريشليه، ونثبت عدم انعدام دوال \(L\) على الخط \(\Re(s)=1\)، ثم نعزل
قطب الشخصية الرئيسية ونطبق مبرهنة Wiener--Ikehara على سلسلة الفئة ذات
المعاملات غير السالبة.

لا يدعي الفصل انتظامًا عندما ينمو \(q\) مع \(x\)، ولا حد خطأ فعالًا، ولا
مبرهنة Siegel--Walfisz، ولا مبرهنة Bombieri--Vinogradov، ولا نتيجة مشروطة
بفرضية ريمان المعممة. تبقى هذه موضوعات مستقلة مؤجلة.

\subsection{المتطلبات السابقة}

يعتمد الفصل على:
\begin{itemize}
  \item الجمع الجزئي واللغة التقاربية من الفصل الثاني.
  \item تعامد شخصيات ديريشليه وبنيتها التحليلية من الفصل السابع.
  \item مرشح الفئة الحسابية من الفصل الثامن.
  \item حد تشيبيشيف، وعدم انعدام زيتا على الخط، وصيغة Wiener--Ikehara من
  الفصل التاسع.
\end{itemize}

\subsection{نواتج التعلم}

بعد إتمام الفصل ينبغي أن يستطيع القارئ:
\begin{enumerate}
  \item تعريف دوال العد الموزونة وغير الموزونة داخل فئة حسابية مختزلة.
  \item تحويل شرط التطابق إلى مجموع على شخصيات ديريشليه.
  \item بناء سلسلة ديريشليه ذات معاملات غير سالبة للفئة نفسها.
  \item إثبات المتراجحة الموزونة لدوال \(L\) حتى عند القيم
  \(|\chi(n)|<1\).
  \item إثبات عدم انعدام \(L(s,\chi)\) على \(\Re(s)=1\) من دون دور.
  \item عزل قطب الشخصية الرئيسية وحساب الباقي \(1/\varphi(q)\).
  \item تطبيق Wiener--Ikehara على \(\psi(x;q,a)\).
  \item الانتقال إلى \(\vartheta(x;q,a)\) و\(\pi(x;q,a)\).
  \item التمييز بين النتيجة النوعية لترديد ثابت والنتائج الكمية الموحدة.
\end{enumerate}

\section{دوال العد في فئة حسابية}

نثبت في بقية الفصل عددًا صحيحًا \(q\ge2\)، وعددًا صحيحًا \(a\) يحقق
\((a,q)=1\). ولكل \(x\ge1\) نعرف
\[
\pi(x;q,a)
=
\sum_{\substack{p\le x\\p\equiv a\pmod q}}1,
\]
\[
\vartheta(x;q,a)
=
\sum_{\substack{p\le x\\p\equiv a\pmod q}}\log p,
\]
و
\[
\psi(x;q,a)
=
\sum_{\substack{n\le x\\n\equiv a\pmod q}}\Lambda(n).
\]

بما أن الفئة \(a\pmod q\) مختزلة، فإن كل عدد \(n\equiv a\pmod q\)
أولي مع \(q\). وهذه الملاحظة تجعل مرشح الشخصيات صالحًا على جميع الأعداد
الداخلة في السلسلة.

\section{سلسلة فون مانغولت في الفئة}

عندما \(\Re(s)>1\)، نضع
\[
F_{q,a}(s)
=
\sum_{\substack{n\ge1\\n\equiv a\pmod q}}
\frac{\Lambda(n)}{n^s}.
\]

\begin{proposition}[تفكيك سلسلة الفئة بواسطة الشخصيات]
\label{prop:class-von-mangoldt-dirichlet-series}
\resultid{ANT-PROP-10-01}
\provedhere
إذا كان \(\Re(s)>1\)، فإن
\[
F_{q,a}(s)
=
\frac1{\varphi(q)}
\sum_{\chi\bmod q}
\overline{\chi(a)}
\left(-\frac{L'}{L}(s,\chi)\right).
\]
\end{proposition}

\begin{proof}
من علاقة التعامد في \textenglish{\texttt{ANT-COR-07-01}}:
\[
\frac1{\varphi(q)}
\sum_{\chi\bmod q}\overline{\chi(a)}\chi(n)
=
\begin{cases}
1,& n\equiv a\pmod q,\\
0,& n\not\equiv a\pmod q.
\end{cases}
\]
إذا كان \((n,q)>1\)، فإن جميع الشخصيات تنعدم عند \(n\)، كما أن
\(n\not\equiv a\pmod q\) لأن \(a\) مختزل؛ لذلك تبقى الهوية صحيحة في هذه
الحالة أيضًا.

وباستعمال التقارب المطلق، ثم
\textenglish{\texttt{ANT-PROP-07-06}}، نحصل على
\begin{align*}
F_{q,a}(s)
&=
\sum_{n=1}^{\infty}
\frac{\Lambda(n)}{n^s}
\frac1{\varphi(q)}
\sum_{\chi\bmod q}\overline{\chi(a)}\chi(n)\\
&=
\frac1{\varphi(q)}
\sum_{\chi\bmod q}\overline{\chi(a)}
\sum_{n=1}^{\infty}\frac{\Lambda(n)\chi(n)}{n^s}\\
&=
\frac1{\varphi(q)}
\sum_{\chi\bmod q}\overline{\chi(a)}
\left(-\frac{L'}{L}(s,\chi)\right).
\end{align*}
\end{proof}

\section{المتراجحة الموزونة لدوال \(L\)}

الحجة الحدية تحتاج موجبية أوسع قليلًا من الهوية المثلثية على دائرة
الوحدة، لأن الشخصية قد تنعدم عند الأعداد غير الأولية مع الترديد.

\begin{lemma}[المتراجحة الموزونة للمشتقات اللوغاريتمية]
\label{lem:weighted-dirichlet-l-log-derivatives}
\resultid{ANT-LEM-10-01}
\provedhere
لكل شخصية ديريشليه \(\chi\pmod q\)، ولكل \(\sigma>1\) و
\(t\in\mathbb R\):
\[
-3\frac{\zeta'}{\zeta}(\sigma)
-4\Re\frac{L'}{L}(\sigma+it,\chi)
-\Re\frac{L'}{L}(\sigma+2it,\chi^2)
\ge0.
\]
\end{lemma}

\begin{proof}
من التقارب المطلق لمتسلسلات المشتقات اللوغاريتمية:
\begin{align*}
&-3\frac{\zeta'}{\zeta}(\sigma)
-4\Re\frac{L'}{L}(\sigma+it,\chi)
-\Re\frac{L'}{L}(\sigma+2it,\chi^2)\\
&\quad=
\sum_{n=1}^{\infty}\frac{\Lambda(n)}{n^\sigma}
\left(3+4\Re z_n+\Re(z_n^2)\right),
\end{align*}
حيث
\[
z_n=\chi(n)n^{-it},
\qquad |z_n|\le1.
\]

يكفي إثبات أن
\[
3+4\Re z+\Re(z^2)\ge0
\qquad(|z|\le1).
\]
اكتب \(z=re^{i\theta}\)، حيث \(0\le r\le1\)، وضع
\(u=\cos\theta\). عندئذ
\[
3+4\Re z+\Re(z^2)
=
2r^2u^2+4ru+3-r^2.
\]
مشتقة الطرف الأيمن بالنسبة إلى \(u\) هي
\[
4r(ru+1)\ge0
\qquad(-1\le u\le1),
\]
لأن \(ru+1\ge1-r\ge0\). لذلك تقع القيمة الدنيا عند \(u=-1\)،
وتساوي
\[
3-4r+r^2=(1-r)(3-r)\ge0.
\]
وبما أن \(\Lambda(n)n^{-\sigma}\ge0\)، تنتج المتراجحة.
\end{proof}

\section{عدم انعدام دوال \(L\) على الخط الواحد}

\begin{theorem}[عدم الانعدام على \(\Re(s)=1\)]
\label{thm:dirichlet-l-nonvanishing-on-one-line}
\resultid{ANT-THM-10-01}
\provedhere
لتكن \(\chi\pmod q\) شخصية ديريشليه.
\begin{enumerate}
  \item إذا كانت \(\chi\) غير رئيسية، فإن
  \[
  L(1+it,\chi)\ne0
  \qquad(t\in\mathbb R).
  \]
  \item إذا كانت \(\chi=\chi_0\) رئيسية، فإن \(L(s,\chi_0)\) لها قطب
  بسيط عند \(s=1\)، وتحقق
  \[
  L(1+it,\chi_0)\ne0
  \qquad(t\ne0).
  \]
\end{enumerate}
\end{theorem}

\begin{proof}
عند \(t=0\)، تعطي
\textenglish{\texttt{ANT-THM-07-09}} عدم انعدام \(L(1,\chi)\) لكل شخصية
غير رئيسية. أما الشخصية الرئيسية فلها قطب بسيط عند \(1\) من
\textenglish{\texttt{ANT-PROP-07-02}}.

افترض الآن أن \(t\ne0\)، وأن \(L(1+it,\chi)\) لها صفر رتبته
\(m\ge1\). عندما \(\sigma\downarrow1\)، يعطينا قطب زيتا البسيط
\[
\frac{\zeta'}{\zeta}(\sigma)
=-\frac1{\sigma-1}+O(1),
\]
ويعطينا التحليل المحلي عند الصفر
\[
\Re\frac{L'}{L}(\sigma+it,\chi)
=
\frac{m}{\sigma-1}+O(1).
\]

لا يمكن أن يكون لـ\(L(s,\chi^2)\) قطب عند \(1+2it\)، لأن قطب الشخصية
الرئيسية الوحيد يقع عند \(1\)، و\(t\ne0\). ولتكن \(m_2\ge0\) رتبة الصفر
عند \(1+2it\)، مع \(m_2=0\) إذا لم تكن النقطة صفرًا. عندئذ
\[
\Re\frac{L'}{L}(\sigma+2it,\chi^2)
=
\frac{m_2}{\sigma-1}+O(1).
\]

إذن الطرف الأيسر في المتراجحة الموزونة يساوي
\[
\frac{3-4m-m_2}{\sigma-1}+O(1).
\]
لكن
\[
3-4m-m_2\le-1,
\]
فيصبح الطرف سالبًا عندما تكون \(\sigma-1\) موجبة وصغيرة، وهذا يناقض
المتراجحة.
\end{proof}

\begin{remark}[الشخصيات غير البدائية]
البرهان السابق يعمل مباشرة لكل شخصية، ولا يحتاج اختزالها أولًا إلى شخصية
بدائية. ويمكن التحقق أيضًا من صيغة الاستحداث: العوامل المحلية المحذوفة
\[
1-\chi^*(p)p^{-s}
\]
لا تنعدم على \(\Re(s)=1\)، لأن
\[
|\chi^*(p)p^{-s}|\le p^{-1}<1.
\]
\end{remark}

\begin{remark}[عدم الدور]
لم يستعمل البرهان مبرهنة الأعداد الأولية في المتتاليات الحسابية، ولا حد
خطأ في الفئات، ولا منطقة خالية كمية، ولا صيغة صريحة للفئات. اعتمد فقط على
المنتجات الأويلرية في \(\Re(s)>1\)، وقطب زيتا، وعدم الانعدام عند
\(s=1\) المثبت في الفصل السابع، والموجبية السابقة.
\end{remark}

\section{عزل قطب الشخصية الرئيسية}

\begin{proposition}[امتداد سلسلة الفئة بعد طرح القطب]
\label{prop:class-series-pole-removal}
\resultid{ANT-PROP-10-02}
\provedhere
تمتد الدالة
\[
F_{q,a}(s)-\frac1{\varphi(q)(s-1)}
\]
هولومورفيًا إلى جوار كل نقطة من الخط \(\Re(s)=1\).
\end{proposition}

\begin{proof}
للشخصية الرئيسية \(\chi_0\pmod q\):
\[
L(s,\chi_0)
=
\zeta(s)P_q(s),
\qquad
P_q(s)=\prod_{p\mid q}(1-p^{-s}).
\]
الدالة \(P_q\) هولومورفية وغير منعدمة في جوار \(s=1\). لذلك
\[
-\frac{L'}{L}(s,\chi_0)
=
-\frac{\zeta'}{\zeta}(s)-\frac{P_q'}{P_q}(s)
=
\frac1{s-1}+H_q(s),
\]
حيث \(H_q\) هولومورفية قرب \(1\).

وبما أن \((a,q)=1\)، فإن \(\chi_0(a)=1\). ومن تفكيك
\textenglish{\texttt{ANT-PROP-10-01}} يكون الباقي عند \(s=1\):
\[
\operatorname*{Res}_{s=1}F_{q,a}(s)
=
\frac1{\varphi(q)}.
\]
جميع حدود الشخصيات غير الرئيسية هولومورفية قرب \(1\)، لأن دوالها تامة
وغير منعدمة عند \(1\) من المبرهنة السابقة.

وعند نقطة \(1+it\) مع \(t\ne0\)، تكون جميع دوال \(L(s,\chi)\)
هولومورفية وغير منعدمة في جوار النقطة؛ فتكون مشتقاتها اللوغاريتمية
هولومورفية. كما أن \(1/(s-1)\) هولومورفية قرب تلك النقطة. ينتج الامتداد
المطلوب على طول الخط كله.
\end{proof}

\section{الصيغة المركزية في الفئة}

نطبق الآن مبرهنة Wiener--Ikehara المسجلة في
\textenglish{\texttt{ANT-THM-09-02}} بحالة \textenglish{\texttt{CITED}}،
ومصدرها \cite[Theorem 1.1، الصفحات 1107--1108]{Korevaar2006}.

\begin{theorem}[مبرهنة الأعداد الأولية الموزونة في فئة ثابتة]
\label{thm:pnt-ap-psi}
\resultid{ANT-THM-10-02}
\provedhere
إذا كان \(q\ge2\) ثابتًا و\((a,q)=1\)، فإن
\[
\psi(x;q,a)
\sim
\frac{x}{\varphi(q)}
\qquad(x\to\infty).
\]
\end{theorem}

\begin{proof}
نطبق Wiener--Ikehara على المعاملات
\[
b_n
=
\Lambda(n)\mathbf 1_{n\equiv a\pmod q}.
\]
هذه المعاملات غير سالبة، وسلسلتها الديريشلية هي \(F_{q,a}(s)\)، المتقاربة
في \(\Re(s)>1\). ومجاميعها الجزئية هي
\[
B(x)=\sum_{n\le x}b_n=\psi(x;q,a).
\]
كما أن
\[
0\le\psi(x;q,a)\le\psi(x)\ll x
\]
من \textenglish{\texttt{ANT-LEM-09-02}}. وأثبتنا أن
\[
F_{q,a}(s)-\frac1{\varphi(q)(s-1)}
\]
يمتد عبر الخط \(\Re(s)=1\). إذن تعطي المبرهنة، مع
\(A=1/\varphi(q)\):
\[
\psi(n;q,a)\sim\frac{n}{\varphi(q)}
\]
للأعداد الصحيحة \(n\). ولعدد حقيقي \(x\):
\[
\psi(x;q,a)=\psi(\lfloor x\rfloor;q,a),
\]
ومن \(\lfloor x\rfloor\sim x\) تنتج الصيغة المطلوبة.
\end{proof}

\section{إزالة القوى الأولية العليا}

\begin{corollary}[الصيغة الموزونة على الأوليات فقط]
\label{cor:pnt-ap-theta}
\resultid{ANT-COR-10-01}
\provedhere
إذا كان \(q\) ثابتًا و\((a,q)=1\)، فإن
\[
\vartheta(x;q,a)
\sim
\frac{x}{\varphi(q)}.
\]
\end{corollary}

\begin{proof}
الفرق \(\psi(x;q,a)-\vartheta(x;q,a)\) هو مساهمة القوى الأولية
\(p^m\le x\) ذات \(m\ge2\) الواقعة في الفئة. لذلك
\[
0\le
\psi(x;q,a)-\vartheta(x;q,a)
\le
\psi(x)-\vartheta(x).
\]
ومن \textenglish{\texttt{ANT-LEM-09-02}}:
\[
\psi(x)-\vartheta(x)
\ll\sqrt{x}\log x=o(x).
\]
وبضم ذلك إلى
\(\psi(x;q,a)=x/\varphi(q)+o(x)\)، نحصل على النتيجة.
\end{proof}

\section{الانتقال إلى دالة العد}

\begin{corollary}[مبرهنة الأعداد الأولية في متتالية حسابية ثابتة]
\label{cor:pnt-ap-pi}
\resultid{ANT-COR-10-02}
\provedhere
إذا كان \(q\ge2\) ثابتًا و\((a,q)=1\)، فإن
\[
\pi(x;q,a)
\sim
\frac{x}{\varphi(q)\log x}.
\]
\end{corollary}

\begin{proof}
يعطي الجمع الجزئي، كما في برهان
\textenglish{\texttt{ANT-THM-02-04}}:
\[
\pi(x;q,a)
=
\frac{\vartheta(x;q,a)}{\log x}
+
\int_2^x
\frac{\vartheta(t;q,a)}{t(\log t)^2}\,dt.
\]
من النتيجة السابقة
\[
\vartheta(t;q,a)
=
\frac{t}{\varphi(q)}+o(t),
\qquad
\vartheta(t;q,a)\ll t.
\]
إذن
\[
\frac{\vartheta(x;q,a)}{\log x}
=
\frac{x}{\varphi(q)\log x}
+o\!\left(\frac{x}{\log x}\right),
\]
بينما
\[
\int_2^x
\frac{\vartheta(t;q,a)}{t(\log t)^2}\,dt
\ll
\int_2^x\frac{dt}{(\log t)^2}
\ll
\frac{x}{(\log x)^2}.
\]
وهذا حد أصغر من الحد الرئيسي، فتنتج الدعوى.
\end{proof}

\section{ماذا تعني النتيجة؟}

لكل فئة مختزلة ثابتة \(a\pmod q\)، يكون نصيبها التقاربي من الكتلة
الأولية مساويًا
\[
\frac1{\varphi(q)}.
\]
أي أن الفئات المختزلة تتقاسم الأوليات بالتساوي في الحد الرئيسي عندما يكون
الترديد ثابتًا. وهذه نتيجة أقوى جوهريًا من مجرد وجود عدد لا نهائي من
الأوليات في كل فئة، المثبت في الفصل الثامن.

لكن عبارة ``بالتساوي'' هنا تقاربية فقط، ولا تعطي سرعة التقارب ولا سلوك
الخطأ لفئة معينة عند ارتفاعات محددة.

\section{ما الذي لم يثبت؟}

لا يثبت المسار الحالي تقديرًا من نوع
\[
\psi(x;q,a)
=
\frac{x}{\varphi(q)}+E(x;q,a)
\]
مع حد فعال أو موحد لـ\(E(x;q,a)\). وللحصول على نتائج كمية نحتاج عادة إلى:
\begin{enumerate}
  \item مناطق خالية كمية لدوال \(L(s,\chi)\).
  \item ضبط الأصفار الحقيقية القريبة من \(1\) وإمكان الصفر الاستثنائي.
  \item صيغ صريحة أو تحويل مسار مع تقديرات موحدة.
  \item أدوات إضافية عندما يتغير \(q\) مع \(x\).
\end{enumerate}

لذلك تبقى مبرهنتا Siegel--Walfisz وBombieri--Vinogradov، والنتائج تحت
GRH، موضوعات منفصلة لاحقة. ولا يجوز نسبتها إلى هذا الفصل.

\section{خريطة الاعتماد المنطقي}

\begin{enumerate}
  \item تعامد الشخصيات يحول شرط التطابق إلى مجموع ملتوي.
  \item المشتقة اللوغاريتمية تحول الالتواء إلى دوال \(L\).
  \item الموجبية الموزونة تمنع أصفار \(L\) على \(\Re(s)=1\).
  \item الشخصية الرئيسية وحدها تعطي قطبًا عند \(1\).
  \item مرشح الفئة يقسم باقي القطب على \(\varphi(q)\).
  \item Wiener--Ikehara تحول القطب وعدم الانعدام إلى
  \(\psi(x;q,a)\sim x/\varphi(q)\).
  \item ضبط القوى الأولية العليا يعطي صيغة \(\vartheta\).
  \item الجمع الجزئي يعطي صيغة \(\pi\).
\end{enumerate}

لا توجد حافة تعود من النتيجة النهائية إلى برهان عدم الانعدام أو إلى
تحليل القطب.

\section{أخطاء شائعة}

\begin{enumerate}
  \item تطبيق Wiener--Ikehara مباشرة على سلسلة ملتوية ذات معاملات مركبة.
  \item استعمال الهوية \(3+4\cos u+\cos2u\ge0\) من دون معالجة حالة
  \(\chi(n)=0\).
  \item الخلط بين عدم انعدام \(L(1,\chi)\) وعدم الانعدام على الخط كله.
  \item إغفال العوامل المحلية عند الشخصيات غير البدائية.
  \item نسيان أن قطب الشخصية الرئيسية وحده هو الذي يولد الحد الرئيسي.
  \item اعتبار النتيجة موحدة في \(q\) مع أنها مثبتة لكل \(q\) ثابت.
  \item استنتاج حد خطأ فعال من صيغة تقاربية تاوبيرية مجردة.
\end{enumerate}

\section{أسئلة تفسير البرهان}

\begin{enumerate}
  \item لماذا نطبق الأداة التاوبيرية على سلسلة الفئة لا على كل التواء منفرد؟
  \item أين استعملت فرضية \((a,q)=1\)؟
  \item لماذا لا يسبب \(L(s,\chi^2)\) قطبًا عند \(1+2it\) عندما
  \(t\ne0\)؟
  \item كيف يظهر العامل \(1/\varphi(q)\) من مرشح الشخصيات؟
  \item ما الفرق المنطقي بين الفصل الثامن والفصل الحالي؟
  \item لماذا لا تكفي النتيجة الحالية لإثبات Siegel--Walfisz؟
\end{enumerate}

\section{تمارين}

\subsection{المستوى الأول}
\begin{enumerate}
  \item تحقق من مرشح الفئة عندما \((n,q)>1\).
  \item أثبت أن \(P_q(1)=\prod_{p\mid q}(1-p^{-1})=\varphi(q)/q\).
  \item بين أن \(\psi(x;q,a)\le\psi(x)\).
  \item أثبت أن فرق \(\psi(x;q,a)-\vartheta(x;q,a)\) غير سالب.
\end{enumerate}

\subsection{المستوى الثاني}
\begin{enumerate}
  \item أثبت مباشرة أن
  \[
  3+4\Re z+\Re(z^2)\ge0
  \qquad(|z|\le1).
  \]
  \item اشتق المتراجحة الموزونة حدًا حدًا من متسلسلات المشتقات
  اللوغاريتمية.
  \item تحقق من أن العوامل المحلية المحذوفة لا تنعدم على \(\Re(s)=1\).
  \item احسب باقي \(F_{q,a}\) عند \(s=1\) من تفكيك الشخصيات.
\end{enumerate}

\subsection{المستوى الثالث}
\begin{enumerate}
  \item طابق فروض Wiener--Ikehara بندًا بندًا مع معاملات سلسلة الفئة.
  \item أعد برهان الانتقال من \(\vartheta(x;q,a)\) إلى \(\pi(x;q,a)\)
  مع تقدير التكامل المتبقي.
  \item حدد الموضع الذي يفشل فيه البرهان إذا سمحنا لـ\(q\) بالنمو مع
  \(x\).
  \item صمم خريطة منفصلة للانتقال من النتيجة النوعية إلى Siegel--Walfisz،
  مع وسم كل أداة غير متاحة حاليًا.
\end{enumerate}

\section{المراجع ومسار التوسع}

للأصل التاريخي انظر \cite{DeLaValleePoussin1896}. وللمعالجات القياسية
للأوليات في المتتاليات الحسابية والخواص التحليلية لدوال \(L\)، راجع
\cite{Davenport2000,MontgomeryVaughan2007,Tenenbaum2015,
IwaniecKowalski2004}.

اعتمد الفصل صيغة Wiener--Ikehara الموثقة في الفصل التاسع من
\cite{Korevaar2006}. أما النتائج الكمية والموحدة فتبقى خارج نطاق هذه
النسخة.

\section{خلاصة الفصل}

حوّل تعامد الشخصيات شرط التطابق إلى مجموع من المشتقات اللوغاريتمية لدوال
\(L\). ومنعت المتراجحة الموزونة أصفار هذه الدوال على الخط \(\Re(s)=1\)،
فبقي قطب الشخصية الرئيسية وحده، وباقيه داخل سلسلة الفئة
\(1/\varphi(q)\). ثم حولت Wiener--Ikehara هذه المعلومة التحليلية إلى
التوزيع المتساوي النوعي للأوليات بين الفئات المختزلة عند الترديد الثابت.
